% ========================================================
% BÀI TẬP TOÁN 5 - KẾT NỐI TRI THỨC VỚI CUỘC SỐNG
% Bài 49: Hình khai triển của hình lập phương, hình hộp chữ nhật, hình trụ
% Định dạng: BTTeX Part III - Trắc nghiệm điền số ngắn
% ========================================================

% Câu 1
\begin{ex}
  Hình khai triển của hình lập phương gồm bao nhiêu mặt hình vuông?
  \shortans{6}
  \loigiai{
    Hình lập phương có 6 mặt hình vuông bằng nhau.
  }
\end{ex}

% Câu 2
\begin{ex}
  Hình lập phương có bao nhiêu đỉnh?
  \shortans{8}
  \loigiai{
    Hình lập phương có 8 đỉnh.
  }
\end{ex}

% Câu 3
\begin{ex}
  Hình lập phương có bao nhiêu cạnh?
  \shortans{12}
  \loigiai{
    Hình lập phương có 12 cạnh.
  }
\end{ex}

% Câu 4
\begin{ex}
  Hình hộp chữ nhật có bao nhiêu mặt?
  \shortans{6}
  \loigiai{
    Hình hộp chữ nhật có 6 mặt.
  }
\end{ex}

% Câu 5
\begin{ex}
  Hình hộp chữ nhật có bao nhiêu đỉnh?
  \shortans{8}
  \loigiai{
    Hình hộp chữ nhật có 8 đỉnh.
  }
\end{ex}

% Câu 6
\begin{ex}
  Hình hộp chữ nhật có bao nhiêu cạnh?
  \shortans{12}
  \loigiai{
    Hình hộp chữ nhật có 12 cạnh.
  }
\end{ex}

% Câu 7
\begin{ex}
  Một hình lập phương có cạnh 6\text{ cm}. Tổng độ dài tất cả các cạnh của nó là bao nhiêu xăng-ti-mét?
  \shortans{72}
  \loigiai{
    Hình lập phương có 12 cạnh bằng nhau. Tổng độ dài: 6 x 12 = 72 cm.
  }
\end{ex}

% Câu 8
\begin{ex}
  Một hình lập phương có cạnh 7\text{ cm}. Tổng độ dài tất cả các cạnh của nó là bao nhiêu xăng-ti-mét?
  \shortans{84}
  \loigiai{
    Hình lập phương có 12 cạnh bằng nhau. Tổng độ dài: 7 x 12 = 84 cm.
  }
\end{ex}

% Câu 9
\begin{ex}
  Một hình lập phương có cạnh 8\text{ cm}. Tổng độ dài tất cả các cạnh của nó là bao nhiêu xăng-ti-mét?
  \shortans{96}
  \loigiai{
    Hình lập phương có 12 cạnh bằng nhau. Tổng độ dài: 8 x 12 = 96 cm.
  }
\end{ex}

% Câu 10
\begin{ex}
  Một hình lập phương có cạnh 9\text{ cm}. Tổng độ dài tất cả các cạnh của nó là bao nhiêu xăng-ti-mét?
  \shortans{108}
  \loigiai{
    Hình lập phương có 12 cạnh bằng nhau. Tổng độ dài: 9 x 12 = 108 cm.
  }
\end{ex}

% Câu 11
\begin{ex}
  Một hình lập phương có cạnh 10\text{ cm}. Tổng độ dài tất cả các cạnh của nó là bao nhiêu xăng-ti-mét?
  \shortans{120}
  \loigiai{
    Hình lập phương có 12 cạnh bằng nhau. Tổng độ dài: 10 x 12 = 120 cm.
  }
\end{ex}

% Câu 12
\begin{ex}
  Một hình lập phương có cạnh 11\text{ cm}. Tổng độ dài tất cả các cạnh của nó là bao nhiêu xăng-ti-mét?
  \shortans{132}
  \loigiai{
    Hình lập phương có 12 cạnh bằng nhau. Tổng độ dài: 11 x 12 = 132 cm.
  }
\end{ex}

% Câu 13
\begin{ex}
  Một hình lập phương có cạnh 12\text{ cm}. Tổng độ dài tất cả các cạnh của nó là bao nhiêu xăng-ti-mét?
  \shortans{144}
  \loigiai{
    Hình lập phương có 12 cạnh bằng nhau. Tổng độ dài: 12 x 12 = 144 cm.
  }
\end{ex}

% Câu 14
\begin{ex}
  Một hình lập phương có cạnh 13\text{ cm}. Tổng độ dài tất cả các cạnh của nó là bao nhiêu xăng-ti-mét?
  \shortans{156}
  \loigiai{
    Hình lập phương có 12 cạnh bằng nhau. Tổng độ dài: 13 x 12 = 156 cm.
  }
\end{ex}

% Câu 15
\begin{ex}
  Một hình lập phương có cạnh 14\text{ cm}. Tổng độ dài tất cả các cạnh của nó là bao nhiêu xăng-ti-mét?
  \shortans{168}
  \loigiai{
    Hình lập phương có 12 cạnh bằng nhau. Tổng độ dài: 14 x 12 = 168 cm.
  }
\end{ex}

% Câu 16
\begin{ex}
  Một hình lập phương có cạnh 15\text{ cm}. Tổng độ dài tất cả các cạnh của nó là bao nhiêu xăng-ti-mét?
  \shortans{180}
  \loigiai{
    Hình lập phương có 12 cạnh bằng nhau. Tổng độ dài: 15 x 12 = 180 cm.
  }
\end{ex}

% Câu 17
\begin{ex}
  Một hình lập phương có cạnh 16\text{ cm}. Tổng độ dài tất cả các cạnh của nó là bao nhiêu xăng-ti-mét?
  \shortans{192}
  \loigiai{
    Hình lập phương có 12 cạnh bằng nhau. Tổng độ dài: 16 x 12 = 192 cm.
  }
\end{ex}

% Câu 18
\begin{ex}
  Một hình lập phương có cạnh 17\text{ cm}. Tổng độ dài tất cả các cạnh của nó là bao nhiêu xăng-ti-mét?
  \shortans{204}
  \loigiai{
    Hình lập phương có 12 cạnh bằng nhau. Tổng độ dài: 17 x 12 = 204 cm.
  }
\end{ex}

% Câu 19
\begin{ex}
  Một hình lập phương có cạnh 18\text{ cm}. Tổng độ dài tất cả các cạnh của nó là bao nhiêu xăng-ti-mét?
  \shortans{216}
  \loigiai{
    Hình lập phương có 12 cạnh bằng nhau. Tổng độ dài: 18 x 12 = 216 cm.
  }
\end{ex}

% Câu 20
\begin{ex}
  Một hình lập phương có cạnh 19\text{ cm}. Tổng độ dài tất cả các cạnh của nó là bao nhiêu xăng-ti-mét?
  \shortans{228}
  \loigiai{
    Hình lập phương có 12 cạnh bằng nhau. Tổng độ dài: 19 x 12 = 228 cm.
  }
\end{ex}

% Câu 21
\begin{ex}
  Một hình lập phương có cạnh 20\text{ cm}. Tổng độ dài tất cả các cạnh của nó là bao nhiêu xăng-ti-mét?
  \shortans{240}
  \loigiai{
    Hình lập phương có 12 cạnh bằng nhau. Tổng độ dài: 20 x 12 = 240 cm.
  }
\end{ex}

% Câu 22
\begin{ex}
  Một hình lập phương có cạnh 21\text{ cm}. Tổng độ dài tất cả các cạnh của nó là bao nhiêu xăng-ti-mét?
  \shortans{252}
  \loigiai{
    Hình lập phương có 12 cạnh bằng nhau. Tổng độ dài: 21 x 12 = 252 cm.
  }
\end{ex}

% Câu 23
\begin{ex}
  Một hình lập phương có cạnh 22\text{ cm}. Tổng độ dài tất cả các cạnh của nó là bao nhiêu xăng-ti-mét?
  \shortans{264}
  \loigiai{
    Hình lập phương có 12 cạnh bằng nhau. Tổng độ dài: 22 x 12 = 264 cm.
  }
\end{ex}

% Câu 24
\begin{ex}
  Một hình lập phương có cạnh 23\text{ cm}. Tổng độ dài tất cả các cạnh của nó là bao nhiêu xăng-ti-mét?
  \shortans{276}
  \loigiai{
    Hình lập phương có 12 cạnh bằng nhau. Tổng độ dài: 23 x 12 = 276 cm.
  }
\end{ex}

% Câu 25
\begin{ex}
  Một hình lập phương có cạnh 24\text{ cm}. Tổng độ dài tất cả các cạnh của nó là bao nhiêu xăng-ti-mét?
  \shortans{288}
  \loigiai{
    Hình lập phương có 12 cạnh bằng nhau. Tổng độ dài: 24 x 12 = 288 cm.
  }
\end{ex}
